%2multibyte Version: 5.50.0.2960 CodePage: 65001
% Master Plan del corso "Metodi Quantitativi per la Finanza"
% Versione aggiornata: nuova Lezione 4 applicativa, sdoppiamento processi,
% eliminazione applicazione autonoma sugli alberi binomiali, introduzione
% del Goal Programming.
% Versione predisposta per Scientific Workplace 5.5.
% Scelte LaTeX conservative: pacchetti standard, tabelle semplici,
% accenti italiani in forma esplicita, ad esempio probabilit\`{a}.

\documentclass[12pt,a4paper]{article}%
\usepackage{amsmath}
\usepackage{amssymb}
\usepackage{array}
\usepackage{longtable}
\usepackage{amsfonts}
\usepackage{graphicx}%
\setcounter{MaxMatrixCols}{30}
%TCIDATA{OutputFilter=latex2.dll}
%TCIDATA{Version=5.50.0.2960}
%TCIDATA{Codepage=65001}
%TCIDATA{LastRevised=Saturday, June 13, 2026 16:15:39}
%TCIDATA{<META NAME="GraphicsSave" CONTENT="32">}
%TCIDATA{<META NAME="SaveForMode" CONTENT="1">}
%TCIDATA{BibliographyScheme=Manual}
%BeginMSIPreambleData
\providecommand{\U}[1]{\protect\rule{.1in}{.1in}}
%EndMSIPreambleData
\setlength{\textwidth}{15.5cm}
\setlength{\textheight}{23.0cm}
\setlength{\oddsidemargin}{0.3cm}
\setlength{\evensidemargin}{0.3cm}
\setlength{\topmargin}{-0.4cm}
\setlength{\parindent}{0pt}
\setlength{\parskip}{6pt}
\newcommand{\R}{\mathbb{R}}
\newcommand{\N}{\mathbb{N}}
\newcommand{\E}{\mathbb{E}}
\newcommand{\Prob}{\mathbb{P}}
\newcommand{\Q}{\mathbb{Q}}
\newcommand{\F}{\mathcal{F}}
\newcommand{\B}{\mathcal{B}}
\newcommand{\G}{\mathcal{G}}
\newcommand{\1}{\mathbf{1}}
\newcommand{\Var}{\operatorname{Var}}
\newcommand{\Cov}{\operatorname{Cov}}
\newcommand{\VaR}{\operatorname{VaR}}
\newcommand{\CVaR}{\operatorname{CVaR}}
\newcommand{\argmin}{\operatorname*{arg\,min}}
\newcommand{\argmax}{\operatorname*{arg\,max}}
\newcommand{\EVPI}{\operatorname{EVPI}}
\newcommand{\VSS}{\operatorname{VSS}}
\begin{document}

\tableofcontents

\begin{center}
{\Large \textbf{Master Plan del corso}}\\[6pt]{\large \textbf{Metodi
Quantitativi per la Finanza Moderna}}\\[12pt]
\end{center}

\textbf{Destinatari:} studenti del quinto anno del corso di laurea in Banca e
Risk Management. Principi di Python.\newline\textbf{Lingua del materiale:}
italiano.\newline\textbf{Formato di scrittura:} LaTeX compatibile con
Scientific Workplace 5.5.\newline\textbf{Prodotti finali:} manuale del corso;
slides delle 16 lezioni; esercizi teorici; applicazioni Python; grafici didattici.

\newpage

\section{Impostazione generale del progetto}

Il corso deve essere progettato come un percorso coerente su modellizzazione
dell'incertezza, valutazione finanziaria e decisione quantitativa sotto
rischio. La sequenza delle lezioni deve evitare l'impressione di una raccolta
eterogenea di strumenti matematici e deve invece costruire progressivamente un
linguaggio comune.

La progressione concettuale aggiornata \`{e} la seguente:

\begin{enumerate}
\item rappresentazione probabilistica dell'incertezza;

\item variabili casuali, distribuzioni e momenti;

\item informazione e valore atteso condizionato;

\item consolidamento computazionale in Python dei concetti probabilistici;

\item processi stocastici in tempo discreto, filtrazioni e martingale;

\item processi stocastici in tempo continuo, diffusioni e simulazione;

\item catene di Markov e rischio di credito;

\item misure di rischio in contesto markoviano;

\item programmazione lineare, con dualit\`{a} come contenuto essenziale;

\item goal programming e decisioni multicriterio;

\item asset allocation multicriterio e asset liability management;

\item programmazione stocastica e decisioni adattate agli scenari.
\end{enumerate}

Il livello didattico desiderato \`{e}: rigoroso nella notazione, selettivo
nelle dimostrazioni, applicativo nell'interpretazione finanziaria. Gli
studenti devono poter seguire il passaggio dalla formulazione matematica al
significato economico-finanziario, con particolare attenzione alla coerenza
dei simboli, alla precisione delle ipotesi e all'uso dei grafici come supporto
alla comprensione.

\section{Criteri generali di scrittura}

\subsection{Manuale}

Il manuale deve essere il riferimento scientifico del corso. Ogni capitolo
dovrebbe contenere:

\begin{enumerate}
\item obiettivi della lezione;

\item motivazione finanziaria;

\item definizioni e notazione;

\item risultati teorici principali;

\item esempi numerici;

\item interpretazioni grafiche;

\item applicazioni finanziarie;

\item esercizi svolti;

\item esercizi proposti;

\item sintesi finale.
\end{enumerate}

Per le lezioni Python si aggiungono:

\begin{enumerate}
\item formulazione matematica del problema computazionale;

\item descrizione dell'algoritmo;

\item codice Python commentato;

\item output numerici;

\item grafici didattici;

\item interpretazione economico-finanziaria dei risultati;

\item esercizi di modifica del codice.
\end{enumerate}

\subsection{Slides}

Le slides devono essere uno strumento di lezione, non una versione compressa
del manuale.

Ogni slide deve essere generata in modo vincolato al Capitolo corrispondente
del Manuale e al Master Plan. \`{E} vietato introdurre definizioni, formule,
propriet\`{a}, esempi, figure, interpretazioni o terminologia tecnica sulla
base della sola conoscenza generale del tema. Un contenuto pu\`{o} entrare
nelle slide solo se \`{e} presente nel Manuale, previsto dal Master Plan o
richiesto esplicitamente dal docente.

Per ogni lezione, la successione delle slide prevede orientativamente:

\begin{enumerate}
\item motivazione e obiettivi;

\item concetti teorici essenziali;

\item esempi numerici guidati;

\item grafici esplicativi;

\item esercizi da svolgere in aula;

\item sintesi conclusiva.
\end{enumerate}

Per le lezioni Python le slides devono privilegiare formulazione, algoritmo,
codice selezionato, output e interpretazione.

\subsection{Scheda operativa: scaletta-promemoria della lezione}

\label{subsec:scheda-scaletta-promemoria}

Per ogni lezione teorica del corso deve essere predisposta una
\emph{scaletta-promemoria} basata sulla versione aggiornata del manuale e
delle slide disponibili nella pagina pubblica di progetto:
\[
\mbox{\texttt{https://paolilla25.synology.me/github-publish/public/}}.
\]
La scaletta non deve essere un riassunto discorsivo della lezione, ma uno
strumento operativo per il docente durante l'erogazione. Deve quindi indicare,
per gruppi omogenei di slide, il contenuto essenziale, il tempo previsto e la
funzione didattica del gruppo.

\paragraph{Struttura obbligatoria di ciascun gruppo.}

Ogni gruppo deve essere presentato secondo il seguente schema:
\[%
\begin{array}
[c]{l}%
\mbox{\texttt{Numero. Titolo del gruppo --- slide x--y}}\\[3pt]%
\mbox{\texttt{Tempo parziale: ... minuti}}\\[3pt]%
\mbox{\texttt{Tempo cumulato: ... minuti}}\\[3pt]%
\mbox{\texttt{Contenuti essenziali: ...}}
\end{array}
\]
Il titolo del gruppo deve essere sintetico e coerente con il titolo o con la
funzione didattica delle slide corrispondenti. Il numero di slide deve essere
indicato con precisione, usando la numerazione effettiva del PDF compilato.

\paragraph{Gestione dei tempi.}

La durata complessiva della lezione teorica \`{e} di $2$ ore e $15$ minuti,
includendo $15$ minuti di pausa. Il tempo effettivo di lezione \`{e} quindi
pari a $120$ minuti. La scaletta deve riportare il tempo parziale di ciascun
gruppo, il tempo cumulato progressivo e il break come voce autonoma.

\paragraph{Contenuti essenziali.}

La voce \emph{Contenuti essenziali} deve essere scritta in forma telegrafica,
non discorsiva. Deve contenere normalmente $3$--$4$ righe, a seconda della
complessit\`{a} del gruppo, e includere le nozioni centrali, le formule
indispensabili, il significato finanziario e le eventuali avvertenze
concettuali importanti.

\paragraph{Uso della notazione matematica.}

Nella scaletta-promemoria prodotta in chat si deve usare la matematica
renderizzata standard della pagina web, non codice LaTeX sorgente. Quando la
scaletta viene invece inserita nel Master Plan in formato LaTeX, le formule
devono essere scritte con notazione LaTeX standard, senza dipendere da comandi
non definiti nell'ambiente di destinazione.

\subsection{Redazione e sintassi matematica}

Nella redazione assistita del manuale, delle slides e dei materiali di
supporto occorre distinguere tra coerenza matematica, coerenza redazionale e
sintassi matematica utilizzata nei diversi contesti di scrittura.

\paragraph{Coerenza matematica.}

Prima di proporre nuovo testo per il manuale, devono essere controllate
esplicitamente le definizioni, le ipotesi, le formule e la coerenza logica dei
passaggi. In particolare, devono essere dichiarate le condizioni necessarie
per la validit\`{a} dei risultati utilizzati: misurabilit\`{a},
integrabilit\`{a}, positivit\`{a} della probabilit\`{a} degli eventi
condizionanti, esistenza di densit\`{a} quando si utilizzano rappresentazioni
continue, e uguaglianze quasi certe quando si lavora con valori attesi
condizionati rispetto a sigma-algebre.

\paragraph{Coerenza redazionale.}

Ogni nuova sezione deve essere scritta come parte di un testo progressivo e
non come unit\`{a} autosufficiente. Prima di introdurre una definizione,
occorre verificare se la stessa nozione sia gi\`{a} stata presentata in
capitoli o sezioni precedenti. Se la nozione \`{e} gi\`{a} stata introdotta,
essa non deve essere ridefinita integralmente, ma richiamata solo nella misura necessaria.

\paragraph{Procedura operativa.}

Per ogni nuova sezione del manuale la redazione assistita deve seguire questa sequenza:

\begin{enumerate}
\item verificare che cosa \`{e} gi\`{a} stato introdotto nei capitoli o nelle
sezioni precedenti;

\item identificare la funzione redazionale della nuova sezione;

\item controllare la correttezza matematica di definizioni, ipotesi, formule e
passaggi logici;

\item verificare la coerenza con il documento di notazione e con il Master Plan;

\item produrre testo LaTeX in blocchi di codice copiabili, evitando formati
intermedi che possano alterare i delimitatori matematici.
\end{enumerate}

\section{Convenzioni preliminari di notazione}

Il documento \texttt{MQF\_Notazione.tex} costituisce il riferimento normativo
per tutti i simboli matematici del corso. Il presente Master Plan riporta
soltanto una sintesi operativa delle convenzioni principali, utile per
collegare ciascuna lezione alla notazione comune.

Ogni lezione deve rispettare tre regole:

\begin{enumerate}
\item non introdurre simboli locali in conflitto con simboli gi\`{a} fissati
in \texttt{MQF\_Notazione.tex};

\item dichiarare esplicitamente eventuali eccezioni locali;

\item aggiornare \texttt{MQF\_Notazione.tex} quando una nuova convenzione
diventa ricorrente nel manuale, nelle slides o nel codice Python.
\end{enumerate}

In coerenza con il documento \texttt{MQF\_Notazione.tex}, si adottano le
seguenti convenzioni generali. Le integrazioni relative al tempo continuo e al
Goal Programming devono essere verificate nel file di notazione prima della
scrittura definitiva dei capitoli corrispondenti.

\begin{center}%
\begin{tabular}
[c]{cp{8.5cm}}\hline
\textbf{Simbolo} & \textbf{Significato}\\\hline
\multicolumn{1}{l}{$(\Omega,\mathcal{F},\mathbb{P})$} & Spazio di
probabilit\`{a}.\\
\multicolumn{1}{l}{$X,Y,Z$} & Variabili casuali reali.\\
\multicolumn{1}{l}{$L$} & Variabile casuale di perdita.\\
\multicolumn{1}{l}{$H$} & Payoff terminale.\\
\multicolumn{1}{l}{$\mathbb{E}[X]$, $\operatorname{Var}(X)$,
$\operatorname{Cov}(X,Y)$} & Valore atteso, varianza e covarianza.\\
\multicolumn{1}{l}{$\{\mathcal{F}_{t}\}_{t=0}^{T}$} & Filtrazione in tempo
discreto.\\
\multicolumn{1}{l}{$\{\mathcal{F}_{t}\}_{t\geq0}$} & Filtrazione in tempo
continuo, quando introdotta.\\
\multicolumn{1}{l}{$X_{t}$} & Processo stocastico generico al tempo $t$.\\
\multicolumn{1}{l}{$S_{t}$} & Prezzo di un'attivit\`{a} finanziaria al tempo
$t$.\\
\multicolumn{1}{l}{$W_{t}$} & Moto browniano standard.\\
\multicolumn{1}{l}{$dW_{t}$} & Incremento browniano infinitesimo nelle SDE.\\
\multicolumn{1}{l}{$\mu$, $\sigma$} & Drift e volatilit\`{a} nei modelli
diffusivi.\\
\multicolumn{1}{l}{$\kappa$, $\theta$} & Velocit\`{a} di mean-reversion e
livello di lungo periodo.\\
\multicolumn{1}{l}{$\rho$, $\Sigma$} & Correlazione e matrice di
covarianza/correlazione nei processi multivariati.\\
\multicolumn{1}{l}{$r$} & Tasso privo di rischio su un periodo, salvo diversa
specificazione.\\
\multicolumn{1}{l}{$r_{t}$} & Tasso istantaneo o fattore di sconto stocastico,
quando introdotto.\\
\multicolumn{1}{l}{$\mathbb{Q}$} & Misura risk-neutral, quando introdotta.\\
\multicolumn{1}{l}{$q$} & Probabilit\`{a} risk-neutral nel modello binomiale,
se usata.\\
\multicolumn{1}{l}{$x$} & Vettore di variabili decisionali nei problemi di
ottimizzazione.\\
\multicolumn{1}{l}{$\lambda$} & Variabile duale o prezzo ombra in
programmazione lineare.\\
\multicolumn{1}{l}{$g_{i}(x)$} & Valore del criterio o obiettivo $i$ in un
modello di Goal Programming.\\
\multicolumn{1}{l}{$b_{i}$, $d_{i}^{+}$, $d_{i}^{-}$} & Target e deviazioni
positive/negative dall'obiettivo $i$.\\
\multicolumn{1}{l}{$\omega\in\Omega$} & Esito elementare nello spazio
campionario.\\
\multicolumn{1}{l}{$\xi$} & Vettore aleatorio o informazione incerta nei
problemi decisionali.\\
\multicolumn{1}{l}{$\mathcal{S}=\{1,\ldots,N\}$} & Insieme finito degli
scenari discreti.\\
\multicolumn{1}{l}{$s\in\mathcal{S}$, $p_{s}$} & Indice di scenario e
probabilit\`{a} dello scenario $s$.\\
\multicolumn{1}{l}{$\xi_{s}$} & Realizzazione dei dati aleatori nello scenario
$s$.\\
\multicolumn{1}{l}{$\mathcal{I}=\{1,\ldots,m\}$} & Insieme finito degli stati
di una catena di Markov.\\
\multicolumn{1}{l}{$M$, $k=1,\ldots,M$} & Numero e indice delle repliche
simulate nelle applicazioni Monte Carlo.\\\hline
\end{tabular}

\end{center}

Il corso utilizza sia strutture in tempo discreto,
\[
t=0,1,\ldots,T,
\]
sia strutture in tempo continuo,
\[
t\geq0,
\]
con una distinzione esplicita tra processi discreti, diffusioni continue e
loro discretizzazione numerica.

\section{Struttura delle 16 lezioni}

\begin{longtable}{>{\centering}p{0.9cm}>{\centering}p{0.8cm}p{4.3cm}p{4.2cm}p{3.3cm}}
\hline
\textbf{Lez.} & \textbf{Tipo} & \textbf{Titolo definitivo} & \textbf{Tema centrale} & \textbf{Output principale} \\
\hline
\endfirsthead
\hline
\textbf{Lez.} & \textbf{Tipo} & \textbf{Titolo definitivo} & \textbf{Tema centrale} & \textbf{Output principale} \\
\hline
\endhead
1 & P & Elementi di teoria della probabilit\`{a} & Spazio di probabilit\`{a}, eventi, probabilit\`{a} condizionata & Capitolo teorico ed esercizi \\
2 & P & Variabili casuali e distribuzioni & Variabili casuali, distribuzioni, momenti & Capitolo teorico, esercizi e grafici \\
3 & P & Valori attesi condizionati e informazione & Condizionamento, partizioni, informazione & Capitolo teorico ed esercizi \\
4 & C & Applicazione in Python: probabilit\`{a}, variabili casuali, valori attesi condizionati & Simulazione, distribuzioni empiriche, stime condizionate & Notebook o script e slides applicative \\
5 & P & Processi stocastici I -- Tempo discreto: fondamenti, informazione e martingale & Processi discreti, filtrazioni, adattabilit\`{a}, martingale & Capitolo teorico ed esercizi \\
6 & P & Processi stocastici II -- Tempo continuo: diffusioni, mean-reversion e simulazione & Browniano, SDE, GBM, OU, CIR, correlazione & Capitolo teorico, esercizi e grafici \\
7 & C & Applicazione in Python: simulazione di traiettorie e pricing & GBM, OU, processi correlati, opzioni asiatiche, OU--CIR & Notebook o script e slides applicative \\
8 & P & Catene di Markov: definizioni e propriet\`{a} & Stati, transizioni, matrice di transizione, distribuzioni stazionarie & Capitolo teorico ed esercizi \\
9 & P & Catene di Markov e misure di rischio & Transizioni di rating, perdite, VaR e CVaR & Capitolo teorico, esercizi e grafici \\
10 & C & Applicazione in Python: rischio di credito & Simulazione di rating e rischio di portafoglio & Notebook o script e slides applicative \\
11 & P & Ottimizzazione e programmazione lineare & Formulazione, soluzioni di base, dualit\`{a} essenziale & Capitolo teorico, esercizi e grafici \\
12 & P & Goal Programming & Obiettivi multipli, deviazioni, priorit\`{a}, trade-off & Capitolo teorico ed esercizi \\
13 & C & Applicazione in Python: Asset Allocation & Asset allocation multicriterio e Asset Liability Management & Notebook o script e slides applicative \\
14 & P & Programmazione lineare stocastica I & Decisioni sotto incertezza, modelli a due stadi, scenari & Capitolo teorico ed esercizi \\
15 & P & Programmazione lineare stocastica II & Modelli multistadio, alberi di scenari, non anticipativit\`{a} & Capitolo teorico ed esercizi \\
16 & C & Applicazione in Python: programmazione stocastica & Asset allocation sotto incertezza e implementazione a scenari & Notebook o script e slides applicative \\
\hline
\end{longtable}


\textbf{Legenda:} P = lezione prevalentemente teorica; C = lezione con forte
componente computazionale o applicativa.

\section{Risorse bibliografiche del corso}

I materiali didattici del corso si basano sui seguenti testi di riferimento.
Per ciascun testo \`{e} indicato il peso relativo attribuito nella redazione
del manuale.

\begin{enumerate}
\item \textbf{Ross, S.M.} \textit{Introduction to Probability Models},
10a~edizione, Elsevier, 2011. \textbf{Peso alto.} Riferimento principale per
processi stocastici, catene di Markov, moto browniano e simulazione.

\item \textbf{Shreve, S.E.} \textit{Stochastic Calculus for Finance~I: The
Binomial Asset Pricing Model}, Springer, 2004. \textbf{Peso medio.} Utilizzato
per filtrazioni in tempo discreto, martingale e modello binomiale come
struttura informativa elementare, non come blocco autonomo di pricing.

\item \textbf{Shreve, S.E.} \textit{Stochastic Calculus for Finance~II:
Continuous-Time Models}, Springer, 2004. \textbf{Peso medio-alto.} Riferimento
per moto browniano, SDE e modelli continui.

\item \textbf{Bj\"{o}rk, T.} \textit{Arbitrage Theory in Continuous Time},
Oxford University Press, 3a~edizione. \textbf{Peso medio.} Utilizzato per il
raccordo tra processi continui e pricing di derivati.

\item \textbf{\O ksendal, B.} \textit{Stochastic Differential Equations: An
Introduction with Applications}, Springer, 6a~edizione. \textbf{Peso medio.}
Riferimento per equazioni differenziali stocastiche e metodo di discretizzazione.

\item \textbf{Hull, J.C.} \textit{Options, Futures, and Other Derivatives},
Pearson, 11a~edizione. \textbf{Peso medio.} Utilizzato per applicazioni
finanziarie, interpretazione dei parametri e pricing Monte Carlo.

\item \textbf{McNeil, A.J., Frey, R., Embrechts, P.} \textit{Quantitative Risk
Management: Concepts, Techniques and Tools}, Princeton University Press,
2a~edizione. \textbf{Peso medio.} Riferimento per misure di rischio, VaR, CVaR
e gestione quantitativa del rischio.

\item \textbf{Glasserman, P.} \textit{Monte Carlo Methods in Financial
Engineering}, Springer, 2003. \textbf{Peso medio.} Riferimento per
simulazione, metodo Euler--Maruyama e stima Monte Carlo.

\item \textbf{Testi di Operations Research e Stochastic Programming da
selezionare.} Riferimenti operativi per programmazione lineare, Goal
Programming, asset allocation multicriterio e programmazione stocastica.
\end{enumerate}

\subsection{Impostazione generale delle lezioni applicative}

Le lezioni applicative hanno la funzione di tradurre in forma computazionale i
principali modelli quantitativi introdotti nelle lezioni teoriche. Esse non
costituiscono un corso autonomo di programmazione, ma un laboratorio di
modellizzazione: il codice Python viene utilizzato come strumento per rendere
operativi, osservabili e verificabili concetti probabilistici, finanziari e di
ottimizzazione gi\`{a} discussi nel manuale e nelle slide teoriche.

Le cinque lezioni applicative del corso sono collocate nelle Lezioni 4, 7, 10,
13 e 16. Esse devono essere sviluppate secondo una struttura comune, in modo che
gli studenti riconoscano una progressione stabile tra problema finanziario,
specificazione teorica, procedura computazionale, output osservabili, controlli
e interpretazione.

Ciascuna lezione applicativa deve prevedere una coppia di casi:
\begin{enumerate}
\item un \textbf{caso aula}, sviluppato dal docente durante la lezione, con uso
esplicito e controllato dell'IA generativa secondo prompt virtuosi;

\item un \textbf{caso take-home}, distinto dal caso aula ma metodologicamente
comparabile, assegnato agli studenti come lavoro autonomo.
\end{enumerate}

Il caso take-home non deve essere una semplice variazione parametrica del caso
aula. Deve invece essere isomorfo sul piano metodologico: deve richiedere
strumenti teorici, tappe logiche e tipi di output analoghi, ma in un contesto
finanziario o probabilistico distinto.

Per ogni lezione applicativa devono essere predisposti almeno i seguenti
materiali:
\begin{enumerate}

\item Traccia del caso aula (\texttt{Caso\_Aula\_Traccia.md}). 
È il caso svolto dal professore. Serve a mostrare il metodo: decomposizione del problema, uso dei regimi A/B/C, prompt virtuosi, costruzione del notebook, controlli e interpretazione.

\item Notebook docente del caso aula (\texttt{Caso\_Aula\_Notebook\_Docente.ipynb}). 
Notebook Jupyter completo, eseguito, con celle Markdown e codice. Non deve essere solo codice: deve mostrare input, output, tappe, prompt esemplari, controlli e interpretazione.

\item Successione di prompt virtuosi utilizzati nel caso aula \\(\texttt{Caso\_Aula\_PromptTrace\_Docente.md}). 
Traccia ordinata dei prompt usati dal docente. Questa non è necessariamente consegnata integralmente agli studenti, ma può essere mostrata in aula come modello.

\item Traccia del caso take-home (\texttt{Caso\_TakeHome\_Traccia.md}). 
Caso diverso ma strutturalmente comparabile al caso aula. Deve richiedere gli stessi strumenti concettuali, ma non consentire una semplice sostituzione meccanica.

\item Scheda docente di calibrazione del caso take-home \\(\texttt{Caso\_TakeHome\_Scheda\_Calibrazione\_Docente.md}). 
Documento interno, non necessariamente distribuito. Deve contenere: premesse teoriche attese, tappe risolutive, input/output tra tappe, output richiesti, prompt docente di riferimento, numero minimo e massimo di prompt.

\item Notebook studente, semi-strutturato (\texttt{Caso\_TakeHome\_Notebook\_Studente.ipynb}). 
Può essere semi-vuoto o strutturato per blocchi. Gli studenti non devono partire da un file bianco, coerentemente con le Guidelines, che raccomandano notebook semi-strutturati con dati, parametri, funzioni, output, grafici, blocchi da completare e domande interpretative.

\item Template del tracciato IA in formato Markdown, che verrà prodotto dallo studente (\texttt{Template\_Tracciato\_AI\_Studente.md}). 
Questo diventa l’oggetto centrale della valutazione metodologica: non solo ``cosa hai ottenuto'', ma ``come hai governato l'AI''.

\item Rubrica di valutazione del notebook e del tracciato IA (\texttt{Rubrica\_Valutazione.md}). 
Documento contenente criteri e pesi di valutazione del notebook e del \\ \texttt{Tracciato\_AI\_Studente}.

\item Prompt da sottoporre all'AI per valutare la qualità del \texttt{Tracciato\_AI\_Studente} (\texttt{Prompt\_Valutatore\_AI.md}). Serve un prompt valutatore che riceva: traccia del caso, scheda docente di calibrazione, notebook/output dello studente e tracciato .md. La valutazione AI deve produrre punteggio, motivazione, prompt migliori, prompt problematici, eventuali deleghe improprie

\end{enumerate}

La scheda docente di calibrazione del caso take-home deve essere predisposta
prima dell'assegnazione agli studenti. Essa deve indicare:
\begin{enumerate}
\item premesse teorico-matematiche necessarie;

\item scomposizione attesa del problema in tappe;

\item collegamenti input--output tra le tappe;

\item output richiesti, quali stime, tabelle, grafici e controlli;

\item successione di prompt docente di riferimento;

\item numero minimo e massimo di prompt ammessi per il tracciato dello studente;

\item criteri di valutazione.
\end{enumerate}

L'uso dell'IA generativa nelle lezioni applicative deve essere organizzato
secondo tre regimi.

\medskip

\noindent
\textbf{Regime A: ricognizione teorico-modellistica.}
Lo studente utilizza l'IA per aiutarsi a riconoscere grandezze
economico-finanziarie, variabili casuali o decisionali, eventi, stati
informativi, scenari, ipotesi, formule e quantit\`{a} teoriche rilevanti.
L'IA non deve risolvere il problema, n\`{e} imporre il modello finale.

\medskip

\noindent
\textbf{Regime B: operazionalizzazione computazionale.}
Data una specifica teorica validata, l'IA pu\`{o} essere utilizzata per
costruire l'apparato computazionale: simulazione, dataset, struttura delle
celle, codice Python, tabelle, output numerici e implementazione tecnica dei
grafici. L'IA non deve modificare variabili, eventi, formule, ipotesi o
significato finanziario del problema.

\medskip

\noindent
\textbf{Regime C: verifica e interpretazione critica.}
Lo studente verifica e interpreta. L'IA pu\`{o} intervenire solo come revisore
critico di controlli, interpretazioni o discussioni gi\`{a} formulate dallo
studente. L'IA non deve produrre l'interpretazione finale al posto dello
studente.

I tre regimi operano su due scale. Al livello del problema complessivo, essi
servono a costruire il percorso risolutivo: inquadramento teorico, tappe
principali, output finali e controlli. Al livello della singola tappa, essi
servono a governare un passaggio specifico della procedura.

Ogni tappa applicativa deve essere trattata come modulo input--output:
\[
\text{input}_{k}
\longrightarrow
\text{operazione}_{k}
\longrightarrow
\text{output}_{k}
\longrightarrow
\text{uso in } k+1.
\]
Per ogni tappa devono essere esplicitati: input disponibili, obiettivo,
regime prevalente, prompt virtuoso di riferimento, oggetti teorici coinvolti,
operazione computazionale, output prodotto, controllo numerico o logico e uso
dell'output nella tappa successiva.

Il notebook Jupyter non deve essere una semplice sequenza di celle di codice.
Ogni tappa deve essere documentata mediante celle Markdown e celle codice,
rendendo visibile il legame tra formulazione teorica, operazione Python, output
e interpretazione. Una tappa didattica pu\`{o} comprendere pi\`{u} celle:
specificazione teorica, prompt di riferimento, codice, controllo e commento
interpretativo.

La valutazione dei lavori take-home deve riguardare congiuntamente il notebook
e il tracciato dell'interazione con l'IA. Non deve premiare la complessit\`{a}
autonoma del codice Python, ma la capacit\`{a} dello studente di governare
l'interazione con l'IA mantenendo controllo teorico, computazionale e
interpretativo del problema quantitativo-finanziario.

La struttura didattica ricorrente delle lezioni applicative \`{e} quindi:
\[
\text{caso aula dimostrativo}
\longrightarrow
\text{caso take-home calibrato}
\longrightarrow
\text{notebook operativo}
\longrightarrow
\text{tracciato IA valutabile}.
\]

\section{Schede dettagliate delle lezioni}

\subsection{Lezione 1 -- Elementi di teoria della probabilit\`{a}}

\textbf{Tipo:} P.\newline\textbf{Domanda guida:} come si rappresenta
formalmente l'incertezza finanziaria?\newline\textbf{Prerequisiti:} algebra
elementare; nozioni intuitive di probabilit\`{a}.\newline\textbf{Obiettivi
didattici:} introdurre spazio campionario, eventi, sigma-algebra,
probabilit\`{a}, probabilit\`{a} condizionata, formula delle probabilit\`{a}
totali e formula di Bayes.\newline\textbf{Notazione introdotta:}
$(\Omega,\mathcal{F},\mathbb{P})$, eventi $A,B$, partizione $\{B_{i}\}$,
probabilit\`{a} condizionata $\mathbb{P}(A\mid B)$.\newline\textbf{Formule
principali:}
\[
\mathbb{P}(A\mid B)=\frac{\mathbb{P}(A\cap B)}{\mathbb{P}(B)},\qquad
\mathbb{P}(B)>0.
\]
\[
\mathbb{P}(A)=\sum_{i} \mathbb{P}(A\mid B_{i})\mathbb{P}(B_{i}),
\qquad\mathbb{P}(B_{j}\mid A)=\frac{\mathbb{P}(A\mid B_{j})\mathbb{P}(B_{j}%
)}{\sum_{i} \mathbb{P}(A\mid B_{i})\mathbb{P}(B_{i})}.
\]
\textbf{Esercizi previsti:} probabilit\`{a} condizionata su eventi di mercato;
aggiornamento bayesiano di uno scenario di default; calcolo di probabilit\`{a}
congiunte.\newline\textbf{Grafici previsti:} diagrammi di Venn; albero
probabilistico a due stadi.\newline\textbf{Applicazioni Python collegate:}
Lezione 4.\newline\textbf{Stato di avanzamento:} bozza progettuale
predisposta; contenuti teorici da sviluppare.

\subsection{Lezione 2 -- Variabili casuali e distribuzioni}

\textbf{Tipo:} P.\newline\textbf{Domanda guida:} come si descrive
quantitativamente un rendimento, una perdita o un payoff aleatorio?\newline%
\textbf{Prerequisiti:} Lezione 1.\newline\textbf{Obiettivi didattici:}
definire variabili casuali discrete e continue; introdurre distribuzioni,
funzione di ripartizione, momenti e quantili; collegare momenti e
rischio.\newline\textbf{Notazione introdotta:} $X$, $F_{X}(x)$, $f_{X}(x)$,
$p_{X}(x)$, $\mathbb{E}[X]$, $\operatorname{Var}(X)$.\newline\textbf{Formule
principali:}
\[
F_{X}(x)=\mathbb{P}(X\leq x), \qquad\mathbb{E}[X]=\sum_{x} x p_{X}(x),
\qquad\mathbb{E}[X]=\int_{-\infty}^{+\infty}x f_{X}(x)\,dx.
\]
\[
\operatorname{Var}(X)=\mathbb{E}[(X-\mathbb{E}[X])^{2}].
\]
\textbf{Esercizi previsti:} payoff discreto; rendimento atteso e varianza;
quantile di una distribuzione discreta di perdite.\newline\textbf{Grafici
previsti:} funzione di massa; densit\`{a}; funzione di ripartizione;
rappresentazione dei quantili.\newline\textbf{Applicazioni Python collegate:}
Lezioni 4, 10 e 13.\newline\textbf{Stato di avanzamento:} bozza progettuale
predisposta; esempi numerici da selezionare.

\subsection{Lezione 3 -- Valori attesi condizionati e informazione}

\textbf{Tipo:} P.\newline\textbf{Domanda guida:} come cambia la valutazione di
un payoff quando arriva nuova informazione?\newline\textbf{Prerequisiti:}
Lezioni 1--2.\newline\textbf{Obiettivi didattici:} introdurre il valore atteso
condizionato rispetto a eventi, distribuzioni condizionate, partizioni e
sigma-algebre informative; interpretare il valore atteso condizionato come
previsione coerente con l'informazione disponibile.\newline\textbf{Notazione
introdotta:} $\mathbb{P}(X=x\mid A)$, $F_{X\mid A}(x)$, $f_{X\mid A}(x)$,
$\mathbb{E}[X\mid A]$, $\mathbb{E}[X\mid\mathcal{G}]$, partizioni finite,
sigma-algebre informative.\newline\textbf{Formule principali:}
\[
\mathbb{E}[X\mid\mathcal{G}]=\sum_{i=1}^{n}\mathbb{E}[X\mid B_{i}%
]\mathbf{1}_{B_{i}}, \qquad\mathcal{G}=\sigma(B_{1},\ldots,B_{n}).
\]
\[
\mathbb{E}[\mathbb{E}[X\mid\mathcal{G}]]=\mathbb{E}[X], \qquad\mathbb{E}%
[\mathbb{E}[X\mid\mathcal{G}_{2}]\mid\mathcal{G}_{1}]=\mathbb{E}%
[X\mid\mathcal{G}_{1}], \qquad\mathcal{G}_{1}\subseteq\mathcal{G}_{2}.
\]
\textbf{Esercizi previsti:} valore atteso condizionato rispetto a eventi;
valore atteso condizionato rispetto a partizioni; propriet\`{a} della
torre.\newline\textbf{Grafici previsti:} albero informativo; rappresentazione
per stati e partizioni; confronto tra distribuzione non condizionata e
condizionata.\newline\textbf{Applicazioni Python collegate:} Lezione 4;
Lezione 16 per la distinzione tra informazione anticipativa e non
anticipativa.\newline\textbf{Stato di avanzamento:} bozza progettuale
predisposta; dimostrazioni essenziali da definire.

\subsection{Lezione 4 -- Applicazione in Python: probabilit\`{a}, variabili
casuali, valori attesi condizionati}


\textbf{Tipo:} C.\newline\textbf{Domanda guida:} come si trasformano
probabilit\`{a}, distribuzioni e valori attesi condizionati in procedure
computabili e simulabili?\newline\textbf{Prerequisiti:} Lezioni 1--3; basi
operative di Python.\newline\textbf{Obiettivi didattici:} simulare variabili
casuali discrete e continue; costruire distribuzioni empiriche, momenti e
quantili; stimare valori attesi e valori attesi condizionati rispetto a eventi
e partizioni; confrontare quantit\`{a} teoriche e stime Monte Carlo.\newline%
\textbf{Notazione introdotta:} nessuna nuova notazione teorica sostanziale;
raccordo computazionale con $X^{(k)}$, $L^{(k)}$, $M$, eventi $A_{g}$, classi
informative e indicatori.\newline\textbf{Formule principali:}
\[
\widehat{\mathbb{E}}[X]=\frac{1}{M}\sum_{k=1}^{M}X^{(k)}, \qquad
\widehat{\mathbb{P}}(A)=\frac{1}{M}\sum_{k=1}^{M}\mathbf{1}_{A}(\omega_{k}).
\]
\[
\widehat{\mathbb{E}}[X\mid A_{g}]=\frac{1}{M_{g}}\sum_{k:\omega_{k}\in A_{g}%
}X^{(k)}, \qquad M_{g}=\#\{k:\omega_{k}\in A_{g}\}.
\]
\textbf{Esercizi previsti:} modificare distribuzioni e parametri; stimare
media, varianza e quantili; calcolare valori attesi condizionati rispetto a
soglie come $X>a$, $L>\ell$ o $S>K$; confrontare partizioni
alternative.\newline\textbf{Grafici previsti:} istogrammi; funzione di
ripartizione empirica; quantili; confronto tra stime condizionate per classi
informative.\newline\textbf{Applicazioni Python collegate:} applicazione
Python 1.\newline\textbf{Stato di avanzamento:} nuova lezione introdotta;
capitolo, slides e codice da scrivere.

\paragraph{Materiali applicativi della lezione.}
Per la Lezione 4 devono essere predisposti: capitolo del manuale, slides
applicative, notebook Python semi-strutturato, script Python completo di
soluzione, figure generate da Python e una breve estensione take-home. Il caso
applicativo guida \`{e} l'analisi simulata di una perdita di portafoglio a un
periodo in presenza di stati informativi di mercato.

\paragraph{Controlli intermedi.}
Ogni tappa computazionale deve concludersi con un controllo numerico o logico
semplice e con una domanda interpretativa. I controlli minimi riguardano:
somma delle probabilit\`{a} degli stati, confronto tra probabilit\`{a} teoriche e
frequenze simulate, assenza di classi vuote, coerenza delle probabilit\`{a} di
superamento soglia, costanza delle medie condizionate entro ciascuno stato e
verifica empirica della relazione
\[
\E\left[\E[L\mid\G]\right]=\E[L].
\]

\paragraph{Estensione take-home.}
L'estensione take-home consiste nel confronto tra una partizione informativa
grossolana e una partizione informativa pi\`{u} fine. Gli studenti devono
costruire le corrispondenti previsioni condizionate della perdita, verificarne
la coerenza con la media non condizionata e commentare l'effetto della maggiore
informazione sulla differenziazione delle previsioni.

\paragraph{Caso applicativo guida.}
Il caso applicativo guida della Lezione 4 \`{e} la simulazione Monte Carlo dei
rendimenti e delle perdite di un portafoglio finanziario a un periodo. Il valore
iniziale del portafoglio \`{e} indicato con $V_{0}>0$, il rendimento aleatorio con
$R$ e la perdita con
\[
L=-V_{0}R.
\]
L'informazione di mercato \`{e} rappresentata da una variabile discreta di stato
$Z\in\{N,V,S\}$, dove $N$ indica un mercato normale, $V$ una fase di volatilit\`{a}
elevata e $S$ uno stato di stress. Gli eventi informativi sono
\[
A_N=\{Z=N\},\qquad A_V=\{Z=V\},\qquad A_S=\{Z=S\},
\]
e generano la sigma-algebra informativa
\[
\G=\sigma(A_N,A_V,A_S).
\]
Condizionatamente allo stato $Z=g$, il rendimento del portafoglio viene simulato
come
\[
R\mid Z=g \sim \mathcal{N}(\mu_g,\sigma_g^2).
\]
Il laboratorio utilizza questo caso per stimare distribuzioni empiriche della
perdita, probabilit\`{a} di superamento di soglie, valori attesi non condizionati,
valori attesi condizionati rispetto agli eventi informativi e la variabile
casuale $\E[L\mid\G]$.

\subsection{Lezione 5 -- Processi stocastici I: tempo discreto, informazione e
martingale}

\textbf{Tipo:} P.\newline\textbf{Domanda guida:} come si rappresenta
l'evoluzione temporale dell'incertezza quando l'informazione cresce nel
tempo?\newline\textbf{Prerequisiti:} Lezioni 1--4.\newline\textbf{Obiettivi
didattici:} introdurre processi stocastici in tempo discreto, traiettorie,
scenari, dipendenza temporale, stazionariet\`{a}, incrementi, filtrazione,
adattabilit\`{a} e martingale; usare il modello binomiale come palestra
strutturale per informazione e dinamica discreta, senza sviluppare una lezione
autonoma di pricing binomiale.\newline\textbf{Notazione introdotta:}
$\{X_{t}\}_{t=0}^{T}$, $\{S_{t}\}_{t=0}^{T}$, $\{\mathcal{F}_{t}\}_{t=0}^{T}$,
traiettorie $X_{0}(\omega),\ldots,X_{T}(\omega)$, scenari discreti
$s\in\mathcal{S}$, dati di scenario $\xi_{s}$.\newline\textbf{Formule
principali:}
\[
X_{t}:\Omega\rightarrow\mathbb{R}, \qquad X_{t} \hbox{ \`{e} } \mathcal{F}%
_{t}\hbox{-misurabile}.
\]
\[
\mathbb{E}[X_{t+1}\mid\mathcal{F}_{t}]=X_{t} \qquad
\hbox{condizione di martingala}.
\]
\[
S_{t+1}=%
\begin{cases}
u S_{t}, & \hbox{stato up},\\
dS_{t}, & \hbox{stato down},
\end{cases}
\qquad\mathcal{F}_{t}%
=\hbox{informazione generata dai primi }t\hbox{ movimenti}.
\]
\textbf{Esercizi previsti:} costruzione di traiettorie discrete;
identificazione della filtrazione; verifica di adattabilit\`{a}; verifica
elementare di una martingala; rappresentazione del binomiale come processo
informativo.\newline\textbf{Grafici previsti:} albero informativo; traiettorie
discrete; ventaglio di scenari; albero binomiale usato come struttura di
filtrazione.\newline\textbf{Applicazioni Python collegate:} Lezione 7;
richiami alla Lezione 16 per scenari e non anticipativit\`{a}.\newline%
\textbf{Stato di avanzamento:} nuova articolazione introdotta; testo da
scrivere a partire dalla precedente scheda sui processi.

\subsection{Lezione 6 -- Processi stocastici II: tempo continuo, diffusioni,
mean-reversion e simulazione}

\textbf{Tipo:} P.\newline\textbf{Domanda guida:} come si modellano fattori
finanziari continui nel tempo e come si preparano alla simulazione
numerica?\newline\textbf{Prerequisiti:} Lezioni 1--5.\newline\textbf{Obiettivi
didattici:} introdurre il moto browniano come limite della passeggiata
aleatoria; presentare equazioni differenziali stocastiche, moto browniano
geometrico, processi mean-reverting Ornstein--Uhlenbeck e CIR, processi
multivariati correlati e discretizzazione di Euler--Maruyama.\newline%
\textbf{Notazione introdotta:} $\{W_{t}\}_{t\geq0}$, $dW_{t}$, $dX_{t}$, drift
$a(t,X_{t})$, diffusione $b(t,X_{t})$, parametri $\mu$, $\sigma$, $\kappa$,
$\theta$, correlazione $\rho$.\newline\textbf{Formule principali:}
\[
W_{t}-W_{s}\sim\mathcal{N}(0,t-s), \qquad0\leq s<t.
\]
\[
dX_{t}=a(t,X_{t})\,dt+b(t,X_{t})\,dW_{t}.
\]
\[
dS_{t}=\mu S_{t}\,dt+\sigma S_{t}\,dW_{t}, \qquad S_{t}=S_{0}\exp\left\{
\left( \mu-\frac{1}{2}\sigma^{2}\right) t+\sigma W_{t}\right\} .
\]
\[
dX_{t}=\kappa(\theta-X_{t})\,dt+\sigma\,dW_{t}, \qquad dr_{t}=\kappa
(\theta-r_{t})\,dt+\sigma\sqrt{r_{t}}\,dW_{t}.
\]
\[
X_{t+\Delta t}\simeq X_{t}+a(t,X_{t})\Delta t+b(t,X_{t})\sqrt{\Delta t}%
\,Z_{t}, \qquad Z_{t}\sim\mathcal{N}(0,1).
\]
\textbf{Esercizi previsti:} momenti del moto browniano; confronto tra GBM e
modello mean-reverting; discretizzazione Euler--Maruyama; ruolo della
correlazione tra fattori.\newline\textbf{Grafici previsti:} traiettorie
browniane; GBM per diversi parametri; traiettorie OU e CIR; confronto tra
processi indipendenti e correlati.\newline\textbf{Applicazioni Python
collegate:} Lezione 7.\newline\textbf{Stato di avanzamento:} nuova lezione
introdotta; completare raccordo con file di notazione.

\subsection{Lezione 7 -- Applicazione in Python: simulazione di traiettorie e
pricing}

\textbf{Tipo:} C.\newline\textbf{Domanda guida:} come si implementano modelli
dinamici simulati e come si stimano prezzi finanziari mediante Monte
Carlo?\newline\textbf{Prerequisiti:} Lezioni 5--6; basi operative di
Python.\newline\textbf{Obiettivi didattici:} simulare traiettorie di GBM, OU,
CIR e processi multivariati correlati; stimare il prezzo di opzioni asiatiche;
implementare una struttura OU--CIR per obbligazioni indicizzate
all'inflazione; valutare prezzi Monte Carlo e incertezza di stima.\newline%
\textbf{Notazione introdotta:} raccordo computazionale con traiettorie
$S_{t}^{(k)}$, fattore inflazione $I_{t}$, tasso istantaneo $r_{t}$, fattore
di sconto $D_{T}^{(k)}$, payoff $H^{(k)}$.\newline\textbf{Formule
principali:}
\[
\widehat{V}_{0}=e^{-rT}\frac{1}{M}\sum_{k=1}^{M}H^{(k)}
\]
per sconto deterministico, oppure
\[
\widehat{V}_{0}=\frac{1}{M}\sum_{k=1}^{M}D_{T}^{(k)}H^{(k)}, \qquad
D_{T}^{(k)}=\exp\left( -\sum_{j=0}^{n-1}r_{t_{j}}^{(k)}\Delta t\right)
\]
per sconto stocastico discretizzato. Per un'opzione asiatica aritmetica,
\[
H^{(k)}=\max\left\{ \frac{1}{n}\sum_{j=1}^{n}S_{t_{j}}^{(k)}-K,0\right\} .
\]
\textbf{Esercizi previsti:} modificare passo temporale e numero di
simulazioni; confrontare GBM e OU; introdurre correlazione tra fattori;
stimare intervalli di errore Monte Carlo; analizzare sensibilit\`{a} rispetto
a strike, volatilit\`{a} e mean-reversion.\newline\textbf{Grafici previsti:}
traiettorie simulate; distribuzione dei payoff; convergenza della media Monte
Carlo; confronto tra prezzi al variare dei parametri.\newline%
\textbf{Applicazioni Python collegate:} applicazione Python 2.\newline%
\textbf{Stato di avanzamento:} nuova lezione introdotta; codice da scrivere.

\subsection{Lezione 8 -- Catene di Markov: definizioni e propriet\`{a}}

\textbf{Tipo:} P.\newline\textbf{Domanda guida:} come si modellano transizioni
probabilistiche tra stati finanziari, ad esempio classi di rating?\newline%
\textbf{Prerequisiti:} Lezioni 1--5.\newline\textbf{Obiettivi didattici:}
definire catene di Markov in tempo discreto; introdurre matrice di
transizione, distribuzioni a uno o pi\`{u} passi, distribuzioni stazionarie e
classificazione degli stati.\newline\textbf{Notazione introdotta:} insieme
degli stati $\mathcal{I}=\{1,\ldots,m\}$, stati $i,j\in\mathcal{I}$, catena
$\{X_{t}\}_{t=0}^{T}$, matrice $P=(p_{ij})$, distribuzione iniziale $\pi_{0}$,
distribuzione $\pi_{t}$.\newline\textbf{Formule principali:}
\[
p_{ij}=\mathbb{P}(X_{t+1}=j\mid X_{t}=i),
\]
\[
\mathbb{P}(X_{t+1}=j\mid X_{t}=i,X_{t-1}=i_{t-1},\ldots,X_{0}=i_{0})=p_{ij}.
\]
\[
\pi_{t+1}=\pi_{t}P, \qquad\pi_{t}=\pi_{0}P^{t}, \qquad\pi=\pi P.
\]
\textbf{Esercizi previsti:} calcolo di transizioni a due periodi;
distribuzione futura di rating; identificazione di stati assorbenti; calcolo
di una distribuzione stazionaria semplice.\newline\textbf{Grafici previsti:}
grafo degli stati; heatmap della matrice di transizione; evoluzione delle
probabilit\`{a} di stato.\newline\textbf{Applicazioni Python collegate:}
Lezione 10.\newline\textbf{Stato di avanzamento:} bozza progettuale
predisposta; rinumerare il capitolo e aggiornare i richiami.

\subsection{Lezione 9 -- Catene di Markov e misure di rischio}

\textbf{Tipo:} P.\newline\textbf{Domanda guida:} come si passa dalla dinamica
degli stati alla misurazione del rischio di perdita?\newline%
\textbf{Prerequisiti:} Lezioni 2 e 8.\newline\textbf{Obiettivi didattici:}
applicare catene di Markov alle transizioni di rating creditizio; costruire
distribuzioni di perdita; calcolare e interpretare VaR e CVaR in un contesto
markoviano.\newline\textbf{Notazione introdotta:} perdita $L$,
$\operatorname{VaR}_{\alpha}(L)$, $\operatorname{CVaR}_{\alpha}(L)$,
esposizioni di portafoglio e stati di default.\newline\textbf{Formule
principali:}
\[
\operatorname{VaR}_{\alpha}(L)=\inf\{\ell\in\mathbb{R}:F_{L}(\ell)\geq
\alpha\}.
\]
\[
\operatorname{CVaR}_{\alpha}(L)=\min_{\eta\in\mathbb{R}}\left\{ \eta+\frac
{1}{1-\alpha}\mathbb{E}[(L-\eta)^{+}]\right\} , \qquad(x)^{+}=\max\{x,0\}.
\]
\textbf{Esercizi previsti:} distribuzione delle perdite da matrice di rating;
calcolo di VaR e CVaR su perdite discrete; confronto tra portafogli con
diverse esposizioni.\newline\textbf{Grafici previsti:} distribuzione delle
perdite; quantile VaR; area di coda CVaR; evoluzione della probabilit\`{a} di
default.\newline\textbf{Applicazioni Python collegate:} Lezione 10.\newline%
\textbf{Stato di avanzamento:} bozza progettuale predisposta; rinumerare il
capitolo e aggiornare i richiami.

\subsection{Lezione 10 -- Applicazione in Python: rischio di credito}

\textbf{Tipo:} C.\newline\textbf{Domanda guida:} come si implementa un modello
di transizione di rating e si stima il rischio di credito di
portafoglio?\newline\textbf{Prerequisiti:} Lezioni 8--9; basi operative di
Python.\newline\textbf{Obiettivi didattici:} implementare una matrice di
transizione; simulare evoluzioni di rating; calcolare probabilit\`{a} di
default, distribuzioni di perdita, VaR e CVaR su portafogli.\newline%
\textbf{Notazione introdotta:} raccordo computazionale con $P$, $\pi_{t}$,
$L$, $\operatorname{VaR}_{\alpha}(L)$ e $\operatorname{CVaR}_{\alpha}%
(L)$.\newline\textbf{Formule principali:}
\[
\pi_{t}=\pi_{0}P^{t}, \qquad\widehat{\mathbb{P}}(\tau\leq T)=\frac{1}{M}%
\sum_{k=1}^{M}\mathbf{1}_{\{\tau^{(k)}\leq T\}}.
\]
\textbf{Esercizi previsti:} modificare la matrice di transizione; introdurre
uno stato assorbente; confrontare portafogli con diverse esposizioni;
analizzare sensibilit\`{a} di VaR e CVaR.\newline\textbf{Grafici previsti:}
heatmap della matrice; probabilit\`{a} di default nel tempo; distribuzione
simulata delle perdite.\newline\textbf{Applicazioni Python collegate:}
applicazione Python 3.\newline\textbf{Stato di avanzamento:} bozza progettuale
predisposta; rinumerare codice e slides.

\subsection{Lezione 11 -- Ottimizzazione e programmazione lineare}

\textbf{Tipo:} P.\newline\textbf{Domanda guida:} come si formula un problema
decisionale vincolato in forma lineare e come si interpreta il valore
marginale dei vincoli?\newline\textbf{Prerequisiti:} algebra lineare
elementare; interpretazione geometrica dei vincoli.\newline\textbf{Obiettivi
didattici:} introdurre variabili decisionali, funzione obiettivo, vincoli
lineari, regione ammissibile, soluzioni di base, valore ottimo e dualit\`{a}
come contenuto essenziale per l'interpretazione economica.\newline%
\textbf{Notazione introdotta:} vettore decisionale $x$, matrice dei vincoli
$A$, vettori $b$ e $c$, valore ottimo $z^{*}$, variabili duali $\lambda
$.\newline\textbf{Formule principali:}
\[
\min_{x} \; c^{\prime}x \qquad\hbox{soggetto a} \qquad Ax\geq b,\quad x\geq0.
\]
\[
\max_{\lambda}\; b^{\prime}\lambda\qquad\hbox{soggetto a} \qquad A^{\prime
}\lambda\leq c,\quad\lambda\geq0.
\]
\textbf{Esercizi previsti:} formulazione di un problema di allocazione;
soluzione grafica in due variabili; identificazione di vincoli attivi;
costruzione elementare del duale e interpretazione dei prezzi ombra.\newline%
\textbf{Grafici previsti:} regione ammissibile; rette di livello; punto
ottimo; interpretazione geometrica di vincoli attivi.\newline%
\textbf{Applicazioni Python collegate:} Lezione 13 e Lezione 16.\newline%
\textbf{Stato di avanzamento:} bozza progettuale predisposta; integrare la
dualit\`{a} nella lezione introduttiva senza conservarla come capitolo autonomo.

\subsection{Lezione 12 -- Goal Programming}

\textbf{Tipo:} P.\newline\textbf{Domanda guida:} come si formalizzano
decisioni finanziarie con pi\`{u} obiettivi, target e trade-off
conflittuali?\newline\textbf{Prerequisiti:} Lezione 11.\newline%
\textbf{Obiettivi didattici:} introdurre la programmazione per obiettivi;
distinguere obiettivi rigidi e target flessibili; rappresentare deviazioni
positive e negative; discutere pesi, priorit\`{a} e trade-off tra criteri
conflittuali.\newline\textbf{Notazione introdotta:} criteri $g_{i}(x)$, target
$b_{i}$, deviazioni $d_{i}^{+}$ e $d_{i}^{-}$, pesi $w_{i}^{+}$ e $w_{i}^{-}%
$.\newline\textbf{Formule principali:}
\[
g_{i}(x)+d_{i}^{-} - d_{i}^{+} = b_{i}, \qquad d_{i}^{+}\geq0, \quad d_{i}%
^{-}\geq0, \quad i=1,\ldots,m.
\]
\[
\min\sum_{i=1}^{m}\left( w_{i}^{+}d_{i}^{+} + w_{i}^{-}d_{i}^{-}\right)
\]
con eventuali priorit\`{a} lessicografiche tra gruppi di obiettivi.\newline%
\textbf{Esercizi previsti:} formulazione di un problema con target di
rendimento e rischio; distinzione tra deviazioni desiderabili e
indesiderabili; confronto tra pesi alternativi; interpretazione dei
trade-off.\newline\textbf{Grafici previsti:} spazio degli obiettivi; frontiera
dei trade-off; confronto tra soluzioni ottenute con pesi diversi.\newline%
\textbf{Applicazioni Python collegate:} Lezione 13.\newline\textbf{Stato di
avanzamento:} nuova lezione introdotta; definire esempio guida coerente con
asset allocation e ALM.

\subsection{Lezione 13 -- Applicazione in Python: Asset Allocation}

\textbf{Tipo:} C.\newline\textbf{Domanda guida:} come si implementa una
decisione di portafoglio multicriterio con obiettivi finanziari e vincoli di
liability matching?\newline\textbf{Prerequisiti:} Lezioni 11--12; basi
operative di Python.\newline\textbf{Obiettivi didattici:} implementare modelli
di asset allocation multicriterio; formulare target di rendimento, rischio,
liquidit\`{a} e liability matching; usare Goal Programming per rappresentare
deviazioni rispetto agli obiettivi; interpretare la sensibilit\`{a} delle
soluzioni rispetto ai pesi.\newline\textbf{Notazione introdotta:} pesi di
portafoglio $x_{i}$, rendimenti attesi $\mu_{i}$, esposizioni $a_{ij}$,
liability target $B_{t}$, deviazioni $d_{i}^{+}$ e $d_{i}^{-}$.\newline%
\textbf{Formule principali:}
\[
\sum_{i} x_{i}=1, \qquad x_{i}\geq0.
\]
\[
\sum_{i} \mu_{i} x_{i}+d_{R}^{-} - d_{R}^{+} = R^{*},
\]
\[
\sum_{i} a_{it}x_{i}+d_{t}^{-} - d_{t}^{+} = B_{t}, \qquad t=1,\ldots,T.
\]
\[
\min\; w_{R}^{-} d_{R}^{-} + w_{R}^{+} d_{R}^{+} + \sum_{t}\left( w_{t}%
^{-}d_{t}^{-}+w_{t}^{+}d_{t}^{+}\right) .
\]
\textbf{Esercizi previsti:} cambiare target di rendimento; modificare pesi
degli obiettivi; imporre vincoli di concentrazione; confrontare soluzioni
orientate a rendimento, rischio e copertura delle passivit\`{a}.\newline%
\textbf{Grafici previsti:} composizione del portafoglio; deviazioni dagli
obiettivi; confronto tra scenari di pesi; profilo asset--liability nel
tempo.\newline\textbf{Applicazioni Python collegate:} applicazione Python
4.\newline\textbf{Stato di avanzamento:} nuova lezione introdotta; codice da scrivere.

\subsection{Lezione 14 -- Programmazione lineare stocastica I}

\textbf{Tipo:} P.\newline\textbf{Domanda guida:} come si formula una decisione
presa prima di conoscere lo scenario futuro?\newline\textbf{Prerequisiti:}
Lezioni 11--13.\newline\textbf{Obiettivi didattici:} introdurre decisioni
sotto incertezza, modelli a due stadi, scenari, valore atteso, decisioni
here-and-now, decisioni di recourse e vincoli di non
anticipativit\`{a}.\newline\textbf{Notazione introdotta:} decisione di primo
stadio $x$, scenario $s\in\mathcal{S}$, dati di scenario $\xi_{s}$, decisione
di secondo stadio $y_{s}$, costo di recourse $d_{s}$.\newline\textbf{Formule
principali:}
\[
\min_{x,y_{s}}\; c^{\prime}x+\sum_{s\in\mathcal{S}}p_{s} d_{s}^{\prime}y_{s}
\]
\[
\hbox{soggetto a} \qquad Ax\geq b, \qquad W_{s} y_{s}\geq h_{s}-T_{s}x, \qquad
y_{s}\geq0, \quad s\in\mathcal{S}.
\]
\textbf{Esercizi previsti:} formulazione di un modello a due stadi; verifica
della non anticipativit\`{a}; interpretazione del recourse; confronto tra
decisione deterministica e decisione stocastica.\newline\textbf{Grafici
previsti:} albero degli scenari a due stadi; struttura primo stadio--secondo
stadio.\newline\textbf{Applicazioni Python collegate:} Lezione 16.\newline%
\textbf{Stato di avanzamento:} bozza progettuale predisposta; aggiornare il
raccordo con la nuova Lezione 13.

\subsection{Lezione 15 -- Programmazione lineare stocastica II}

\textbf{Tipo:} P.\newline\textbf{Domanda guida:} come si estende una decisione
stocastica a pi\`{u} date informative e come si impedisce l'uso anticipato
dell'informazione futura?\newline\textbf{Prerequisiti:} Lezione 14; richiami
su filtrazioni dalla Lezione 5.\newline\textbf{Obiettivi didattici:}
introdurre modelli multistadio, alberi di scenari, nodi informativi, decisioni
adattate alla storia osservata e vincoli di non anticipativit\`{a}; mantenere
EVPI e VSS come indicatori sintetici del valore dell'informazione e della
soluzione stocastica.\newline\textbf{Notazione introdotta:} nodi $n$
dell'albero, predecessore $a(n)$, probabilit\`{a} di nodo $p_{n}$, decisione
$x_{n}$, gruppi di scenari con storia comune, $z^{SP}$, $z^{WS}$, $z^{EV}$,
$\operatorname{EVPI}$, $\operatorname{VSS}$.\newline\textbf{Formule
principali:}
\[
x_{s,t}=x_{s^{\prime},t} \qquad\hbox{se gli scenari }s\hbox{ e }s^{\prime
}\hbox{ hanno la stessa storia fino a }t.
\]
Per problemi formulati in minimizzazione,
\[
z^{WS}\leq z^{SP}\leq z^{EV}, \qquad\operatorname{EVPI}=z^{SP}-z^{WS},
\qquad\operatorname{VSS}=z^{EV}-z^{SP}.
\]
\textbf{Esercizi previsti:} costruzione di un piccolo albero multistadio;
identificazione di vincoli di non anticipativit\`{a}; calcolo di EVPI e VSS su
esempio ridotto.\newline\textbf{Grafici previsti:} albero multistadio; gruppi
di scenari con storia comune; confronto dei valori SP, EV e WS.\newline%
\textbf{Applicazioni Python collegate:} Lezione 16.\newline\textbf{Stato di
avanzamento:} bozza progettuale predisposta; riformulare il capitolo mettendo
al centro multistadio e non anticipativit\`{a}.

\subsection{Lezione 16 -- Applicazione in Python: programmazione stocastica}

\textbf{Tipo:} C.\newline\textbf{Domanda guida:} come si implementa un modello
di asset allocation sotto incertezza con scenari e decisioni non
anticipative?\newline\textbf{Prerequisiti:} Lezioni 14--15; basi operative di
Python.\newline\textbf{Obiettivi didattici:} implementare modelli di
programmazione stocastica per problemi di asset allocation sotto incertezza;
costruire scenari; formulare decisioni di primo stadio e recourse; imporre
vincoli di non anticipativit\`{a}; calcolare soluzioni SP, EV e WS; stimare
EVPI e VSS.\newline\textbf{Notazione introdotta:} raccordo computazionale con
$x$, $y_{s}$, $p_{s}$, $\xi_{s}$, $z^{SP}$, $z^{WS}$, $z^{EV}$,
$\operatorname{EVPI}$ e $\operatorname{VSS}$.\newline\textbf{Formule
principali:}
\[
\operatorname{EVPI}=z^{SP}-z^{WS}, \qquad\operatorname{VSS}=z^{EV}-z^{SP}.
\]
\textbf{Esercizi previsti:} modificare probabilit\`{a} degli scenari;
modificare rendimenti e passivit\`{a}; confrontare soluzioni SP, EV e WS;
introdurre o rimuovere vincoli di non anticipativit\`{a}.\newline%
\textbf{Grafici previsti:} albero degli scenari; composizione del portafoglio
per scenario o nodo; confronto dei valori; visualizzazione di EVPI e
VSS.\newline\textbf{Applicazioni Python collegate:} applicazione Python
5.\newline\textbf{Stato di avanzamento:} bozza progettuale predisposta; codice
da scrivere.

\section{Registro degli esercizi}

\begin{longtable}{>{\centering}p{1.1cm}p{4.1cm}p{4.1cm}p{4.5cm}}
\hline
\textbf{Lez.} & \textbf{Esercizio 1} & \textbf{Esercizio 2} & \textbf{Esercizio grafico/computazionale}\\
\hline
\endfirsthead
\hline
\textbf{Lez.} & \textbf{Esercizio 1} & \textbf{Esercizio 2} & \textbf{Esercizio grafico/computazionale}\\
\hline
\endhead
1 & Probabilit\`{a} condizionata & Formula di Bayes & Albero probabilistico\\
2 & Media e varianza di payoff discreto & Quantile di perdita & Grafico della funzione di ripartizione\\
3 & Valore atteso condizionato su partizione & Propriet\`{a} della torre & Albero informativo\\
4 & Simulazione di variabili casuali & Stima condizionata su eventi & Istogrammi, quantili e partizioni\\
5 & Filtrazione e adattabilit\`{a} & Verifica di martingala & Albero binomiale come struttura informativa\\
6 & Moto browniano e GBM & OU, CIR e mean-reversion & Traiettorie simulate e discretizzazione\\
7 & Pricing Monte Carlo di asiatica & Sistema OU--CIR & Grafico traiettorie e payoff\\
8 & Transizioni a due periodi & Stato assorbente e stazionariet\`{a} & Grafo di Markov\\
9 & VaR discreto in modello markoviano & CVaR discreto & Distribuzione delle perdite\\
10 & Transizioni di rating simulate & Probabilit\`{a} di default & Heatmap e distribuzione perdite\\
11 & Formulazione PL & Duale essenziale e prezzi ombra & Regione ammissibile\\
12 & Formulazione Goal Programming & Trade-off tra target & Spazio degli obiettivi\\
13 & Asset allocation multicriterio & Liability matching & Portafogli e deviazioni dai target\\
14 & Modello a due stadi & Recourse e valore atteso & Albero degli scenari a due stadi\\
15 & Non anticipativit\`{a} multistadio & EVPI e VSS & Albero multistadio e confronto valori\\
16 & Modifica degli scenari & Confronto SP/EV/WS & Grafico dei risultati\\
\hline
\end{longtable}


\section{Registro delle applicazioni Python}

\begin{longtable}{>{\centering}p{1.5cm}>{\centering}p{1.4cm}p{4.4cm}p{6.5cm}}
\hline
\textbf{Applic.} & \textbf{Lez.} & \textbf{Tema} & \textbf{Output computazionale} \\
\hline
\endfirsthead
\hline
\textbf{Applic.} & \textbf{Lez.} & \textbf{Tema} & \textbf{Output computazionale} \\
\hline
\endhead
1 & 4 & Probabilit\`{a}, variabili casuali e valori attesi condizionati & Simulazioni di variabili casuali, distribuzioni empiriche, momenti, quantili, stime condizionate. \\
2 & 7 & Simulazione di traiettorie e pricing & Traiettorie GBM, OU, CIR e correlate; opzioni asiatiche; obbligazioni indicizzate all'inflazione con struttura OU--CIR. \\
3 & 10 & Rischio di credito & Matrici di transizione, evoluzioni di rating, probabilit\`{a} di default, VaR e CVaR di portafoglio. \\
4 & 13 & Asset Allocation & Goal Programming, asset allocation multicriterio, asset liability management, analisi dei trade-off. \\
5 & 16 & Programmazione stocastica & Scenari, recourse, vincoli di non anticipativit\`{a}, soluzioni SP/EV/WS, EVPI e VSS. \\
\hline
\end{longtable}


Criteri comuni per il codice Python:

\begin{enumerate}
\item codice ben commentato;

\item separazione tra dati, parametri, funzioni e output;

\item uso di nomi coerenti con la notazione matematica;

\item grafici leggibili;

\item interpretazione dei risultati dopo ogni blocco computazionale;

\item esercizi di modifica del codice.
\end{enumerate}

\section{File previsti del progetto}

\subsection{Manuale}
\begin{verbatim}
/manuale
  MQF_Manuale_Master.tex
/Capitoli 
MQF_Cap_01_Probabilita.tex 
MQF_Cap_02_Variabili_Casuali.tex 
MQF_Cap_03_Valori_Attesi_Condizionati.tex 
MQF_Cap_04_Python_Probabilita_Condizionamento.tex 
MQF_Cap_05_Processi_Stocastici_Tempo_Discreto.tex 
MQF_Cap_06_Processi_Stocastici_Tempo_Continuo.tex 
MQF_Cap_07_Python_Advanced_Pricing.tex 
MQF_Cap_08_Catene_Markov.tex 
MQF_Cap_09_Markov_Misure_Rischio.tex 
MQF_Cap_10_Python_Rischio_Credito.tex 
MQF_Cap_11_Programmazione_Lineare.tex 
MQF_Cap_12_Goal_Programming.tex 
MQF_Cap_13_Python_Asset_Allocation_ALM.tex 
MQF_Cap_14_Programmazione_Stocastica_I.tex 
MQF_Cap_15_Programmazione_Stocastica_II.tex 
MQF_Cap_16_Python_Programmazione_Stocastica.tex
\end{verbatim}

\subsection{Slides}
\begin{verbatim}
/slides
Slides_Lez_01_Probabilita.tex 
Slides_Lez_02_Variabili_Casuali.tex 
Slides_Lez_03_Valori_Attesi_Condizionati.tex 
Slides_Lez_04_Python_Probabilita_Condizionamento.tex 
Slides_Lez_05_Processi_Stocastici_Tempo_Discreto.tex 
Slides_Lez_06_Processi_Stocastici_Tempo_Continuo.tex 
Slides_Lez_07_Python_Advanced_Pricing.tex 
Slides_Lez_08_Catene_Markov.tex 
Slides_Lez_09_Markov_Misure_Rischio.tex 
Slides_Lez_10_Python_Rischio_Credito.tex 
Slides_Lez_11_Programmazione_Lineare.tex 
Slides_Lez_12_Goal_Programming.tex 
Slides_Lez_13_Python_Asset_Allocation_ALM.tex 
Slides_Lez_14_Programmazione_Stocastica_I.tex 
Slides_Lez_15_Programmazione_Stocastica_II.tex 
Slides_Lez_16_Python_Programmazione_Stocastica.tex
\end{verbatim}

\subsection{Python}
\begin{verbatim}
/python
  MQF_Python_04_Probabilita_Variabili_Condizionamento.py
  MQF_Python_07_Python_Advanced_Pricing.py
  MQF_Python_10_Rischio_Credito.py
  MQF_Python_13_Asset_Allocation_Goal_Programming.py
  MQF_Python_16_Programmazione_Stocastica.py
\end{verbatim}

\subsection{Documenti di coordinamento}
\begin{verbatim}
/progetto
  MQF_Master_Plan.tex
  MQF_Notazione.tex
  MQF_Project_Guidelines.md
  MQF_Stato_Avanzamento.md
  MQF_Registro_Esercizi.tex
  MQF_Registro_Grafici.tex
  MQF_Registro_Decisioni.tex
\end{verbatim}

\section{Stato di avanzamento generale}

\begin{longtable}{p{4.2cm}p{3.0cm}p{7.0cm}}
\hline
\textbf{Area} & \textbf{Stato} & \textbf{Note} \\
\hline
\endfirsthead
\hline
\textbf{Area} & \textbf{Stato} & \textbf{Note} \\
\hline
\endhead
Struttura delle 16 lezioni & Aggiornata & Introdotta nuova Lezione 4 applicativa; processi stocastici sdoppiati; eliminata applicazione autonoma sugli alberi binomiali; introdotto Goal Programming. \\
Notazione generale & Bozza consolidata, da integrare & Documento \texttt{MQF\_Notazione.tex} predisposto; verificare integrazioni su tempo continuo, SDE, OU/CIR, correlazione e Goal Programming. \\
Manuale & In sviluppo & Capitoli 1--3 gi\`{a} avviati; capitoli 4--16 da riallineare alla nuova architettura. \\
Slides & In sviluppo & Le slide devono seguire la nuova numerazione; eliminare riferimenti alla vecchia lezione Python sugli alberi binomiali. \\
Esercizi teorici & Bozza aggiornata & Registro esercizi riallineato alla nuova sequenza delle lezioni. \\
Applicazioni Python & Bozza aggiornata & Cinque applicazioni previste: Lezioni 4, 7, 10, 13 e 16. \\
Grafici & Bozza da validare & Aggiornare registro grafici in base alle nuove lezioni su processi, Goal Programming e asset allocation. \\
Compatibilit\`{a} SWP 5.5 & In verifica & Preambolo mantenuto conservativo; evitare pacchetti LaTeX moderni non necessari nei file principali. \\
\hline
\end{longtable}


\section{Questioni aperte}

\begin{enumerate}
\item Verificare e, se necessario, integrare \texttt{MQF\_Notazione.tex} per
tempo continuo, SDE, processi OU/CIR, correlazione e Goal Programming.

\item Definire il titolo definitivo del Capitolo 4 applicativo.

\item Definire il livello di dettaglio del modello binomiale nella Lezione 5,
mantenendolo come esempio strutturale e non come blocco autonomo di pricing.

\item Selezionare l'esempio guida della Lezione 7 per obbligazioni indicizzate
all'inflazione con sistema OU--CIR.

\item Selezionare l'esempio guida della Lezione 13 per asset allocation
multicriterio e Asset Liability Management.

\item Definire il livello computazionale delle librerie di ottimizzazione
Python da usare nelle Lezioni 13 e 16.

\item Aggiornare le Guidelines di progetto in coerenza con la nuova architettura.
\end{enumerate}

\section{Prossimi passi operativi}

\begin{enumerate}
\item Aggiornare \texttt{MQF\_Project\_Guidelines.md} recependo la nuova
sequenza delle 16 lezioni.

\item Aggiornare \texttt{MQF\_Notazione.tex} limitatamente ai simboli che
diventano ricorrenti nei nuovi capitoli.

\item Rinominare o predisporre i file sorgente dei capitoli 4--16 secondo la
nuova numerazione.

\item Predisporre la struttura del Capitolo 4 e della prima applicazione Python.

\item Riallineare progressivamente manuale, slides, esercizi e registri alla
nuova sequenza.
\end{enumerate}


\end{document}